\documentclass[11pt,a4paper]{article}

% ───────────────────────────── Packages ─────────────────────────────
\usepackage[utf8]{inputenc}
\usepackage[T1]{fontenc}
\usepackage{amsmath, amssymb, amsthm}
\usepackage{physics}          % \ket, \bra, \braket
\usepackage{graphicx}
\usepackage{booktabs}
\usepackage{longtable}
\usepackage{array}
\usepackage{geometry}
\usepackage{xcolor}
\usepackage{listings}
\usepackage{enumitem}
\usepackage{hyperref}
\usepackage{caption}
\usepackage{subcaption}
\usepackage{todonotes}

\geometry{margin=2.5cm}

\hypersetup{
    colorlinks=true,
    linkcolor=blue!50!black,
    citecolor=blue!50!black,
    urlcolor=blue!50!black
}

% ───────────────────────────── Code listing style ─────────────────────────────
\definecolor{codebg}{rgb}{0.96,0.96,0.96}
\lstdefinestyle{mystyle}{
    backgroundcolor=\color{codebg},
    basicstyle=\ttfamily\small,
    breaklines=true,
    frame=single,
    showstringspaces=false,
    columns=fullflexible,
    keepspaces=true
}
\lstset{style=mystyle}

% ───────────────────────────── Title ─────────────────────────────
\title{\textbf{Q-MODE: Amino Acid Lattice Mapping for\\ Binding-Pocket Representation and \\ Quantum-Assisted Docking Site Search}}
\author{Author Name(s) \\ \small Multidisciplinary Project}
\date{\today}

\begin{document}

\maketitle

\begin{abstract}
Q-MODE is a computational pipeline that converts a protein binding pocket, supplied as a PDB structure, into a discrete, spatially-consistent sequence of pharmacophoric interaction sites, and subsequently encodes that sequence into qubit-ready quantum states. The pipeline combines classical cheminformatics (RDKit-based feature extraction, Crippen and Abraham solvation-derived intensities, solvent-accessibility filtering, covalent-topology-based ordering) with a quantum-inspired encoding and search stage, following the two-step scheme (first/second encoding) and the modified Grover-search formulation introduced in \emph{``Quantum algorithm for protein-ligand docking sites identification in the interaction space''}. This report documents the full functioning of the pipeline as currently implemented: the residue- and ligand-parsing stage, the pharmacophore and intensity-assignment stage, the surface and deduplication filters, the topological chain assembly, the quantum encodings, and the Grover-search-based candidate ranking. The validation and benchmarking section is intentionally left incomplete in this revision and will be completed in a subsequent iteration of the report.
\end{abstract}

\tableofcontents
\newpage

% ═══════════════════════════════════════════════════════════════════
\section{Introduction}
% ═══════════════════════════════════════════════════════════════════

Protein-ligand docking site identification is a central problem in structure-based drug design: given a protein binding pocket, one wants to identify which sub-regions of that pocket are chemically compatible with a candidate ligand's pharmacophoric profile. Classical docking pipelines typically rely on scoring functions evaluated over a continuous 3D search space, which is expensive to sample exhaustively.

Q-MODE explores an alternative representation: it reduces the binding pocket to a \emph{flat, ordered sequence of pharmacophoric sites} — a ``1D chain'' analogous in spirit to a residue sequence — and then encodes short, ligand-sized windows of that chain into low-qubit-count quantum states. This representation is directly compatible with the amplitude- and basis-state-encoding schemes used by quantum search algorithms (in particular, a modified Grover search), enabling a quantum-inspired (and, in the encoding stage, genuinely quantum-circuit-executed) approach to scanning candidate docking windows.

\subsection{Motivation}

Two properties of a binding pocket motivate the lattice/chain representation used here:

\begin{itemize}[leftmargin=1.5em]
    \item \textbf{Locality}: a ligand interacts with a contiguous, spatially compact set of pocket residues. Representing the pocket as an ordered sequence lets a ``sliding window'' of fixed size stand in for a candidate docking sub-site.
    \item \textbf{Discretization}: encoding each site's chemical character (hydrophobic vs. polar/charged) as a small number of qubits per site keeps the resulting quantum circuits within the qubit counts realistically simulable today (or executable on near-term hardware), while still preserving the essential pharmacophoric signal needed to discriminate candidate sites.
\end{itemize}

\subsection{Scientific Background}

The quantum-encoding and quantum-search stages of Q-MODE follow the scheme proposed in \emph{``Quantum algorithm for protein-ligand docking sites identification in the interaction space''} (Liliopoulos et al.):

\begin{itemize}[leftmargin=1.5em]
    \item A \textbf{first encoding} that binarizes, per site, a hydrophobicity channel ($h$) and a hydrogen-bonding/charge channel ($hb$) into a 2-qubit basis state, so that a window of $k$ sites maps onto a $2k$-qubit computational basis state. This is the representation used to build the Grover oracle.
    \item A \textbf{second encoding} that maps the same two channels onto two pairs of probability amplitudes $(a,b)$ and $(c,d)$, suitable for amplitude-based (e.g. SWAP-test) similarity/distance estimation between two sites or windows.
    \item A \textbf{modified Grover search}, run per ``protein shift'' (i.e. per tiling offset of the pocket chain into non-overlapping ligand-sized windows), that identifies which windows' first-encoding bitstring matches the ligand's own bitstring with probability above the uniform ($1/N$) threshold.
\end{itemize}

Q-MODE implements the first encoding, the second encoding, and the modified Grover search (oracle and diffusion operator, executed on a Qiskit simulator backend) as described above. The paper's own SWAP-test-based ranking of candidate sites is \textbf{not yet implemented}; Q-MODE currently substitutes a classical Euclidean-distance-to-ligand-centroid ranking for benchmark validation, as detailed in Section~\ref{sec:validation}.

% ═══════════════════════════════════════════════════════════════════
\section{Pipeline Overview}
% ═══════════════════════════════════════════════════════════════════

Given a pocket PDB file (and, optionally, a full PDB structure containing a bound ligand), Q-MODE executes the following stages in sequence:

\begin{enumerate}[leftmargin=1.8em]
    \item \textbf{Residue extraction} — parse residues and real 3D atomic coordinates from the PDB file; assign bond orders by structural template comparison rather than positional overlay (Section~\ref{sec:residue-extraction}).
    \item \textbf{Pharmacophoric feature computation} — RDKit-based atom/group feature extraction (Section~\ref{sec:feature-extraction}).
    \item \textbf{Intensity assignment} — Crippen LogP contributions for hydrophobic sites; Abraham solute descriptors for H-bond donor/acceptor/ionizable sites (Section~\ref{sec:intensity}).
    \item \textbf{Surface filtering} — SASA-based removal of solvent-inaccessible (buried) sites (Section~\ref{sec:surface-filter}).
    \item \textbf{Site deduplication} — merging of near-duplicate same-type sites into a single centroid site (Section~\ref{sec:dedup}).
    \item \textbf{Topological ordering} — within each residue, ordering sites by a breadth-first traversal of the covalent bond graph, starting from the backbone nitrogen (Section~\ref{sec:topological-order}).
    \item \textbf{Chain assembly} — concatenation of residues in protein-chain order into one flat sequence (Section~\ref{sec:chain-assembly}).
    \item \textbf{Quantum encoding} — first (binary) and second (amplitude) encodings of ligand-sized sliding-window segments (Section~\ref{sec:quantum-encoding}).
    \item \textbf{Ligand extraction} \emph{(optional)} — parsing of a bound ligand's HETATM group and construction of its own pharmacophore chain (Section~\ref{sec:ligand-extraction}).
    \item \textbf{Grover search} \emph{(optional)} — quantum-circuit-based identification of candidate docking windows matching the ligand's encoded profile (Section~\ref{sec:grover-search}).
    \item \textbf{Distance-to-ligand validation} \emph{(optional, benchmark mode)} — see Section~\ref{sec:validation}, left to be completed.
\end{enumerate}

Figure~\ref{fig:pipeline-flow} sketches the data flow between stages.

\begin{figure}[h]
    \centering
    \fbox{\parbox{0.9\textwidth}{\centering
    \texttt{PDB file} $\rightarrow$ \texttt{ResidueRecord[]} $\rightarrow$ \texttt{AtomFeature[]} (per residue)
    $\rightarrow$ intensity assignment $\rightarrow$ surface filter $\rightarrow$ dedup $\rightarrow$ topological order
    $\rightarrow$ \textbf{flat\_chain} $\rightarrow$ quantum encoding (first/second) $\rightarrow$ \textbf{quantum\_chain}
    \\[0.3em]
    (optional) \texttt{ligand PDB} $\rightarrow$ ligand chain $\rightarrow$ Grover search against \textbf{flat\_chain}
    $\rightarrow$ \textbf{grover\_search} candidates $\rightarrow$ distance-to-ligand validation
    }}
    \caption{High-level data flow of the Q-MODE pipeline.}
    \label{fig:pipeline-flow}
\end{figure}

% ═══════════════════════════════════════════════════════════════════
\section{Methods}
% ═══════════════════════════════════════════════════════════════════

\subsection{Residue Extraction and Bond-Order Assignment}
\label{sec:residue-extraction}

Each residue is read directly from the PDB's \texttt{ATOM}/\texttt{HETATM} records (\texttt{pdb\_reader.py}), keeping only the primary alternate location (\texttt{altLoc} blank or \texttt{A}) and, by default, discarding water. Residues are grouped by $(\texttt{chain\_id}, \texttt{res\_seq}, \texttt{res\_name})$ and sorted by $(\texttt{chain\_id}, \texttt{res\_seq})$ to preserve chain order.

For each residue, an RDKit molecule is built from the real heavy-atom coordinates via a PDB block, and bond orders/aromaticity are assigned by \emph{template comparison} (\texttt{AssignBondOrdersFromTemplate}) against a canonical SMILES template for each of the 20 standard amino acids (plus histidine protonation variants \texttt{HSD/HSE/HSP/HIE/HID/HIP} and the disulfide-bonded \texttt{CYX}). Two templates are used per residue type: a ``free amino acid'' template (full \texttt{-COOH}) when a terminal \texttt{OXT} atom is present, and an ``internal'' template (backbone \texttt{-OH} removed) otherwise — since most residues in a chain are peptide-bonded and not free termini. Missing hydrogens (not resolved by X-ray) are added afterward with 3D coordinates consistent with the heavy-atom geometry (\texttt{AddHs(addCoords=True)}).

This template-based approach is deliberately preferred over a purely geometric/positional bond-order guess: it is robust to the moderate coordinate noise present in crystallographic structures, at the cost of requiring the residue type to be one of the recognized standard (or near-standard) amino acids.

\subsection{Pharmacophoric Feature Extraction}
\label{sec:feature-extraction}

Each residue's RDKit molecule is passed through RDKit's chemical feature factory (\texttt{BaseFeatures.fdef}), which identifies six pharmacophore families, mapped onto the internal type vocabulary used throughout Q-MODE:

\begin{center}
\begin{tabular}{ll}
\toprule
\textbf{RDKit family} & \textbf{Q-MODE type} \\
\midrule
Donor & \texttt{HBondDonor} \\
Acceptor & \texttt{HBondAcceptor} \\
Hydrophobe / LumpedHydrophobe & \texttt{Hydrophobe} \\
Aromatic & \texttt{Aromatic} \\
PosIonizable & \texttt{PosIonizable} \\
NegIonizable & \texttt{NegIonizable} \\
\bottomrule
\end{tabular}
\end{center}

Each feature carries the 3D centroid of its constituent atom(s) and, for multi-atom groups, an atom-index list used by later stages (intensity re-assignment, surface filtering, topological mapping). Where the feature factory is unavailable, a manual fallback identifies donors/acceptors from N/O/F valence and hydrogen count, and hydrophobic sites from per-atom Crippen contributions directly.

\subsection{Intensity Assignment}
\label{sec:intensity}

Every site carries a scalar \emph{intensity}, intended to reflect the physical strength of that site's interaction potential rather than a placeholder constant.

\paragraph{Hydrophobic sites.} \texttt{Hydrophobe} sites are assigned the sum of positive per-atom Crippen LogP contributions (\texttt{rdMolDescriptors.\_CalcCrippenContribs}) over their constituent atoms, falling back to a neutral value of $1.0$ when no positive contribution is present.

\paragraph{H-bond donor/acceptor/ionizable sites.} These sites are assigned an \emph{intrinsic}, per-site hydrogen-bonding intensity from Abraham's Linear Solvation Energy Relationship (LSER) scales — specifically $\alpha_2^H$ (donor strength) for \texttt{HBondDonor}/\texttt{PosIonizable} sites and $\beta_2^H$ (acceptor strength) for \texttt{HBondAcceptor}/\texttt{NegIonizable} sites (\texttt{abraham\_hbond.py}). Values are looked up in a table keyed by $(\texttt{res\_name}, \texttt{pdb\_atom\_name})$ (with a wildcard residue name for backbone atoms shared by all amino acids), cross-checked against the UFZ-LSER database. When the exact residue/atom key is not present — most notably for ligand atoms, which never match a standard-amino-acid key — a heuristic functional-group classifier (\texttt{\_heuristic\_group\_key}) inspects the atom's local chemical environment (element, valence, neighboring atoms) and reuses the closest protein-side proxy already in the table (carboxyl, amine, phenol, thiol); anything outside those four groups keeps the neutral default intensity.

Two known approximations are carried through explicitly rather than silently:
\begin{itemize}[leftmargin=1.5em]
    \item Abraham's scales are defined for \emph{neutral} solutes. Ionized side chains (protonated Lys/Arg, deprotonated Asp/Glu) are proxied by their nearest neutral analog (e.g. propylamine for the lysine ammonium group), which is expected to \emph{underestimate} the true interaction strength of the charged form. Each affected table entry is flagged with an explicit \texttt{note} field.
    \item The ligand's own H-bond intensity has no direct Abraham entry (the table is indexed by amino-acid residue/atom names), so it is covered only by the heuristic group-key fallback described above, and otherwise remains at the neutral default.
\end{itemize}

\subsection{Surface Filtering}
\label{sec:surface-filter}

Buried sites are not physically reachable by an incoming ligand and would otherwise add noise to the pocket representation. Q-MODE filters sites by per-atom Solvent-Accessible Surface Area (SASA), computed with BioPython's Shrake--Rupley implementation (probe radius $1.4$~\r{A}) over the \emph{whole} structure (\texttt{surface\_filter.py}). A site is kept if at least one of its constituent atoms has $\text{SASA} > \tau$, with default threshold $\tau = 1.0$~\r{A}\textsuperscript{2} (a conservative value; the literature range for ``significantly exposed'' spans roughly 1.0--5.0~\r{A}\textsuperscript{2}). This filter is on by default and can be disabled with \texttt{--no-surface-filter}, e.g. to retain buried sites for structural comparison rather than ligandability analysis.

Because RDKit atom indices and PDB atom-name bookkeeping can become misaligned during template-based bond assignment, the matching between a feature and its SASA value is done by \emph{nearest 3D coordinate} to a known PDB atom rather than by index, which is more robust to that misalignment.

\subsection{Site Deduplication (Centroid Merging)}
\label{sec:dedup}

Within a residue, several nearby atoms of the same pharmacophore type (e.g. multiple ring carbons contributing separate \texttt{Aromatic} features) would otherwise appear as separate, highly-correlated sites. \texttt{site\_dedup\_centroid.py} greedily groups same-type features whose pairwise Euclidean distance is below $1.5$~\r{A}, replacing each group with a single site at the group's coordinate centroid and an aggregate intensity equal to the \emph{sum} of the group's individual intensities. Same-type sites that are spatially far apart (e.g. two separate H-bond acceptors on opposite sides of a residue) are correctly kept as distinct sites. The result is additionally capped to the \texttt{max\_sites} most intense sites per residue (default 12 for protein residues; configurable and set lower, e.g. 3, for the ligand side to keep the Grover-search qubit count tractable).

\subsection{Topological Ordering}
\label{sec:topological-order}

Within a residue, the deduplicated sites are ordered by a breadth-first search (BFS) over the residue's covalent bond graph, starting from the (lowest-index) backbone nitrogen atom (\texttt{site\_selection.py}). Each site is first mapped to its nearest heavy atom by 3D distance, after which the BFS visit order over the bond graph determines the site's position in the residue-local sequence. This produces a chemically meaningful, deterministic ordering (ties broken by sorted neighbor traversal) that reflects the residue's actual connectivity rather than an arbitrary or purely spatial order.

\subsection{Chain Assembly}
\label{sec:chain-assembly}

Residues are concatenated in protein-chain order — i.e. by $(\texttt{chain\_id}, \texttt{res\_seq})$ as already established at the parsing stage — and each residue's internally-ordered sites are appended in sequence. The result is one flat list, \texttt{flat\_chain}, where every entry carries: the originating residue label, residue name/sequence/chain id, 3D coordinates, pharmacophore type, and (channel-appropriate) intensity. This flat chain is the shared representation consumed by both the visualization utilities and the quantum-encoding stage.

\subsection{Quantum Encoding}
\label{sec:quantum-encoding}

Each site in \texttt{flat\_chain} is reduced to two scalar channels (\texttt{qubit\_chain.get\_h\_hb\_intensities}):
\[
(h, hb) = \begin{cases}
(\text{intensity}, 0) & \text{type} \in \{\texttt{Hydrophobe}, \texttt{Aromatic}\} \\
(0, \text{intensity}) & \text{type} \in \{\texttt{HBondDonor}, \texttt{HBondAcceptor}, \texttt{PosIonizable}, \texttt{NegIonizable}\}
\end{cases}
\]

\paragraph{First encoding (basis-state binarization).} Each channel is binarized against a per-channel threshold, producing a 2-bit basis state per site:
\[
q_h = \begin{cases} 0 & h < h_{\text{thr}} \\ 1 & h \ge h_{\text{thr}} \end{cases}
\qquad
q_{hb} = \begin{cases} 0 & hb < hb_{\text{thr}} \\ 1 & hb \ge hb_{\text{thr}} \end{cases}
\]
The paper's reference scheme thresholds at the channel's range midpoint $(v_{\min}+v_{\max})/2$. Empirically, real pocket intensity distributions are right-skewed (most sites sit near the floor, with a few strong outliers), which drags the mean upward and leaves almost nothing above threshold — collapsing the chain to an all-zero bitstring. Q-MODE instead thresholds at the \emph{median} of each channel's active (nonzero) values, which splits the population roughly 50/50 and keeps the binarization informative for the downstream Grover search (\texttt{compute\_h\_hb\_thresholds}).

A window of $k$ consecutive sites is thus mapped onto a $2k$-bit basis state $\ket{x} = \ket{q_h^{(1)} q_{hb}^{(1)} \ldots q_h^{(k)} q_{hb}^{(k)}}$.

\paragraph{Second encoding (amplitude encoding).} Each channel value $v$ is also mapped onto a pair of normalized probability amplitudes via a cosine/sine-like projection of $v$ within its channel's range $[v_{\min}, v_{\max}]$:
\[
p = v_{\max} - v, \qquad q = v - v_{\min}, \qquad
(a, b) = \left( \frac{p}{\sqrt{p^2+q^2}}, \; \frac{q}{\sqrt{p^2+q^2}} \right)
\]
applied independently to $h$ (giving $a,b$ with $a^2+b^2=1$) and to $hb$ (giving $c,d$ with $c^2+d^2=1$), with $v$ clamped to $[v_{\min}, v_{\max}]$ beforehand so that a site outside its own channel's range (e.g. $h=0$ for a purely polar site) lands at the range's lower extreme instead of producing an invalid (unnormalizable) state. This representation is intended for amplitude-based (e.g. SWAP-test) similarity/distance estimation between sites, as prescribed by the reference paper.

\paragraph{Sliding-window segmentation.} \texttt{qubit\_chain.build\_qubit\_chain} slides a window of size \texttt{ligand\_size} (default 3 sites) over \texttt{flat\_chain} with unit stride, producing one \texttt{QubitSegment} per window position, each carrying its first-encoding bitstring and its per-site second-encoding amplitudes. This is the ``classical'' quantum-encoding step (Step 7 in the pipeline enumeration), independent of the real ligand used later by Grover search.

\subsection{Ligand Extraction}
\label{sec:ligand-extraction}

When a full PDB structure containing a bound ligand is supplied (\texttt{--ligand-pdb}), \texttt{ligand\_reader.py} parses its \texttt{HETATM} records, excluding water and the 20 standard amino acids, and groups atoms by $(\texttt{chain\_id}, \texttt{res\_seq}, \texttt{res\_name})$ — i.e. by physical copy, since a small/symmetric structure's asymmetric unit may contain several bound copies of the same ligand. The largest non-water, non-amino-acid group (by heavy-atom count) is auto-selected by default, or a specific one via \texttt{--ligand-code}.

Bond orders for the ligand are assigned, in order of preference, by:
\begin{enumerate}[leftmargin=1.8em]
    \item Fetching the ideal SMILES for the ligand's 3-letter code from the RCSB PDB Chemical Component Dictionary (CCD) REST API and applying the same template-based assignment used for amino acids (\texttt{AssignBondOrdersFromTemplate}).
    \item Falling back to geometric bond-order perception (\texttt{rdDetermineBonds}) when the CCD entry is unavailable (unknown code, or no network access).
\end{enumerate}

The extracted ligand is then passed through the same local pipeline as a protein residue — feature extraction, Abraham/heuristic intensity assignment, centroid deduplication (capped to \texttt{--ligand-max-sites}, default 3), and topological ordering — yielding a ligand-side chain of the same shape as a protein window, suitable for direct comparison in the Grover search stage.

\subsection{Grover Search}
\label{sec:grover-search}

For a ligand chain of $k$ sites, Q-MODE runs a modified Grover search (\texttt{qmode/grover/search.py}) against every non-overlapping tiling of the protein's \texttt{flat\_chain} into windows of size $k$:

\begin{enumerate}[leftmargin=1.8em]
    \item \textbf{Tiling.} For each shift offset $o \in \{0, \ldots, k-1\}$, the pocket chain is partitioned into non-overlapping windows of size $k$ starting at $o$ (\texttt{tile\_offset}). Each window is reduced to its first-encoding bitstring; when several windows collapse onto the same bitstring, the one with the highest aggregate $h/hb$ intensity across its sites (the ``most interactive'' window) is retained as that bitstring's representative.
    \item \textbf{Superposition state.} The unique bitstrings observed at this offset are combined into an equal superposition,
    \[
    \ket{s} = \frac{1}{\sqrt{N}} \sum_{i=1}^{N} \ket{x_i},
    \]
    over $n_{\text{qubits}} = 2k$ qubits, where $N$ is the number of unique windows at this offset.
    \item \textbf{Oracle.} The ligand's own first-encoding bitstring $x_{\text{lig}}$ defines a phase-flip oracle,
    \[
    \hat{O} = I - 2\ket{x_{\text{lig}}}\bra{x_{\text{lig}}}.
    \]
    \item \textbf{Diffusion operator.} The Grover diffusion operator about the superposition state,
    \[
    \hat{G} = 2\ket{s}\bra{s} - I,
    \]
    is applied after the oracle.
    \item \textbf{Circuit execution.} $\ket{s}$ is prepared with Qiskit's \texttt{StatePreparation}, followed by the oracle and diffusion operators (each compiled from their dense unitary matrix via \texttt{UnitaryGate}), then measured. The circuit is transpiled and executed on Qiskit Aer's \texttt{AerSimulator} (default $4096$ shots), yielding an empirical probability for each measured bitstring.
    \item \textbf{Acceptance threshold.} A shift offset's search is considered a match if the ligand bitstring's measured probability exceeds the uniform baseline $1/N$ — i.e. the oracle+diffusion step measurably amplified that state above what a uniform superposition alone would give.
\end{enumerate}

Across all shift offsets, matching windows are collected into a candidate list, each carrying: the shift offset, the window's start index in \texttt{flat\_chain}, its \emph{interactivity score} (aggregate $h/hb$ intensity), the residues it spans, the matched bitstring, its measured probability, the $1/N$ threshold, and the number of unique states at that offset. Candidates are sorted by descending interactivity score, so that the most chemically interactive matching window is reported first.

\paragraph{Scalability.} Oracle and diffusion operators are synthesized from dense $2^{n_{\text{qubits}}} \times 2^{n_{\text{qubits}}}$ unitary matrices. Qiskit's generic unitary-synthesis routine does not scale gracefully past roughly 10 qubits (tens of seconds to minutes per shift offset in practice), which is why \texttt{--ligand-max-sites} defaults to 3 (6 qubits); raising it substantially increases circuit-compilation time.

\subsection{Distance-to-Ligand Validation}
\label{sec:validation}

\todo[inline, color=yellow!40]{Validation and benchmarking section — to be completed. This includes: the classical Euclidean-distance-to-ligand-centroid stand-in currently implemented in \texttt{qmode/grover/evaluate.py} (as an interim substitute for the reference paper's SWAP-test-based ranking); quantitative results over the expanded benchmark set (Astex Diverse Set targets plus the four hand-picked complexes in \texttt{make\_pockets.py}); accuracy/success-rate metrics; and a discussion of where the current ranking succeeds or fails relative to the true bound pose.}

% ═══════════════════════════════════════════════════════════════════
\section{Implementation}
% ═══════════════════════════════════════════════════════════════════

\subsection{Repository Structure}

\begin{lstlisting}[language=bash]
Q-MODE-project/
|-- qmode/
|   |-- pdb_reader.py           # PDB parsing -> residues, real 3D coords
|   |-- ligand_reader.py        # Ligand HETATM parsing (CCD template / geometric fallback)
|   |-- feature_extraction.py   # RDKit pharmacophore features + Crippen h
|   |-- abraham_hbond.py        # Abraham hb intensity lookup
|   |-- surface_filter.py       # SASA-based solvent-exposure filter
|   |-- site_selection.py       # Topological (BFS) site ordering
|   |-- site_dedup_centroid.py  # Near-duplicate site merging
|   |-- quantum_encoding.py     # First/second quantum encoding
|   |-- qubit_chain.py          # Sliding-window segmentation + qubit chain
|   |-- grover/
|   |   |-- search.py           # Oracle, diffusion, circuit execution
|   |   `-- evaluate.py         # Distance-to-ligand validation
|   |-- lattice_fitting.py      # 3D->2D PCA projection (utility, unused by main pipeline)
|   |-- snapping.py             # 2D->integer lattice nodes (utility, unused)
|   `-- labeling.py             # One-hot pharmacophore labeling (utility, unused)
|-- scripts/
|   |-- run_pipeline.py         # Main CLI entry point
|   |-- plot_residue_chain.py   # Single-residue 3D + 1D visualization
|   `-- make_pockets.py         # Downloads sample PDBs, crops binding pockets
|-- data/raw/                   # Input / generated pocket PDB files
`-- tests/                      # Unit tests
\end{lstlisting}

\subsection{Dependencies}

\begin{center}
\begin{tabular}{ll}
\toprule
\textbf{Package} & \textbf{Minimum version} \\
\midrule
rdkit & 2023.3.1 \\
numpy & 1.24 \\
pandas & 2.0 \\
scikit-learn & 1.3 \\
scipy & 1.11 \\
matplotlib & 3.7 \\
biopython & 1.79 \\
qiskit & 2.2 \\
qiskit-aer & 0.17 \\
tqdm & 4.65 \\
\bottomrule
\end{tabular}
\end{center}

\subsection{Command-Line Interface}

The pipeline is exposed through a single script, \texttt{scripts/run\_pipeline.py}:

\begin{lstlisting}[language=bash]
python scripts/run_pipeline.py --pdb data/raw/3PTB_pocket.pdb --plot

python scripts/run_pipeline.py --pdb data/raw/3PTB_pocket.pdb \
    --output data/processed/ --save-plot data/processed/pocket.png

python scripts/run_pipeline.py --pdb data/raw/4dfr_pocket.pdb \
    --ligand-pdb data/raw/4DFR.pdb
\end{lstlisting}

\begin{center}
\begin{tabular}{@{}p{3.2cm}p{2cm}p{8.5cm}@{}}
\toprule
\textbf{Option} & \textbf{Default} & \textbf{Description} \\
\midrule
\texttt{--pdb} & required & Path to the pocket (or whole-protein) PDB file. \\
\texttt{--output} & \texttt{None} & Directory for JSON/CSV output files. \\
\texttt{--plot} & \texttt{False} & Show an interactive visualization. \\
\texttt{--save-plot} & \texttt{None} & Save plots as PNG images. \\
\texttt{--no-surface-filter} & filter on & Disable the SASA solvent-exposure filter. \\
\texttt{--sasa-threshold} & \texttt{1.0} & SASA threshold (\r{A}\textsuperscript{2}) for solvent exposure. \\
\texttt{--ligand-size} & \texttt{3} & Sliding-window size (sites) for the classical quantum-chain segmentation (Section~\ref{sec:quantum-encoding}). \\
\texttt{--ligand-pdb} & \texttt{None} & Path to a PDB file with the ligand's HETATM records; enables Grover search. \\
\texttt{--ligand-code} & \texttt{None} & 3-letter ligand code to disambiguate multiple HETATM groups. \\
\texttt{--ligand-max-sites} & \texttt{3} & Max ligand pharmacophore sites used by Grover search (qubit-count control). \\
\bottomrule
\end{tabular}
\end{center}

\subsection{Output Artifacts}

Running the pipeline with \texttt{--output} produces up to three JSON/CSV artifacts:

\paragraph{\texttt{pocket\_chain.json} / \texttt{.csv}} The full flat lattice chain with per-site metadata (residue label, coordinates, pharmacophore type, intensity).

\paragraph{\texttt{quantum\_chain.json}} Sliding-window segments with their first-encoding bitstring and second-encoding amplitudes.

\paragraph{\texttt{grover\_search.json}} \emph{(only when \texttt{--ligand-pdb} is given)} Grover candidates ranked by descending interactivity score, each including its distance-to-ligand validation field once Section~\ref{sec:validation} is completed.

\subsection{Testing}

Unit tests (\texttt{tests/}) cover the ligand reader, the Grover search circuit construction, and an end-to-end pipeline smoke test:

\begin{lstlisting}[language=bash]
pytest tests/
\end{lstlisting}

% ═══════════════════════════════════════════════════════════════════
\section{Known Limitations}
% ═══════════════════════════════════════════════════════════════════

\begin{itemize}[leftmargin=1.5em]
    \item \textbf{Abraham proxies for ionized groups.} Abraham's LSER scales are defined for neutral solutes; ionized side chains (protonated Lys/Arg, deprotonated Asp/Glu) are approximated by their nearest neutral analog, which likely underestimates true interaction strength.
    \item \textbf{Ligand H-bond intensity coverage.} The Abraham table is indexed by amino-acid residue/atom name and does not natively cover ligand atoms; only four functional-group classes (carboxyl, amine, phenol, thiol) are covered by the heuristic fallback, with everything else defaulting to a neutral intensity.
    \item \textbf{Qubit-count scalability.} Qiskit's generic unitary synthesis for the Grover oracle/diffusion operators does not scale past roughly 10 qubits, bounding \texttt{--ligand-max-sites} to small values (default 3, i.e. 6 qubits) for practical runtimes.
    \item \textbf{SWAP-test ranking not implemented.} The reference paper's own quantum (amplitude/SWAP-test-based) ranking of candidate docking sites is not implemented; Q-MODE substitutes a classical Euclidean-distance-to-ligand-centroid ranking, valid only in benchmark mode where the true ligand pose is known (see Section~\ref{sec:validation}, to be completed).
    \item \textbf{Utility modules not wired into the main pipeline.} \texttt{lattice\_fitting.py}, \texttt{snapping.py}, and \texttt{labeling.py} implement a 3D$\to$2D PCA projection, integer-lattice snapping, and one-hot pharmacophore labeling respectively, but are not currently called from \texttt{run\_pipeline.py}.
\end{itemize}

% ═══════════════════════════════════════════════════════════════════
\section{Conclusion and Future Work}
% ═══════════════════════════════════════════════════════════════════

Q-MODE implements a complete, end-to-end path from a raw PDB binding pocket to a quantum-circuit-executed Grover search over candidate docking windows, grounded in physically-derived (Crippen/Abraham) site intensities rather than placeholder values. The classical feature-extraction and chain-assembly stages (Sections~\ref{sec:residue-extraction}--\ref{sec:chain-assembly}) are stable and unit-tested; the quantum-encoding and Grover-search stages (Sections~\ref{sec:quantum-encoding}--\ref{sec:grover-search}) faithfully reproduce the reference paper's first encoding, second encoding, oracle, and diffusion operator.

The immediate next step for this project is completing Section~\ref{sec:validation}: a quantitative validation of the Grover-search candidate ranking against the expanded benchmark set (the Astex Diverse Set targets plus the originally hand-picked complexes), reporting distance-to-ligand accuracy and discussing whether — and how — a genuine SWAP-test-based amplitude ranking would change the results relative to the current classical distance stand-in.

% ═══════════════════════════════════════════════════════════════════
\section*{References}
% ═══════════════════════════════════════════════════════════════════

\begin{enumerate}
    \item Liliopoulos et al., \emph{Quantum algorithm for protein-ligand docking sites identification in the interaction space}.
    \item Abraham, M. H. \emph{Scales of solute hydrogen-bonding: their construction and application to physicochemical and biochemical processes}, Chem. Soc. Rev. (1993/1994) and subsequent updates (2000--2012); values cross-checked against the UFZ-LSER Database (\url{https://www.ufz.de/lserd}).
    \item Wildman, S. A.; Crippen, G. M. \emph{Prediction of Physicochemical Parameters by Atomic Contributions}, J. Chem. Inf. Comput. Sci. (1999).
    \item Shrake, A.; Rupley, J. A. \emph{Environment and exposure to solvent of protein atoms. Lysozyme and insulin}, J. Mol. Biol. (1973).
    \item Hartshorn, M. J. et al. \emph{Diverse, High-Quality Test Set for the Validation of Protein-Ligand Docking Performance}, J. Med. Chem. (2007) — source of the Astex Diverse Set benchmark targets.
\end{enumerate}

\end{document}
